% GENERATED FILE. Edit canonical lessons, then run npm run generate:books.
% canonical-source-hash: fnv1a64:b3cad3cc4ddc10ae
% canonical-lessons: UR-C416-kura, UR-C416-nali, UR-C416-gatar, UR-C416-badbu, UR-C416-phailana

\chapter{Rubbish and Drains}
\label{ch:ur-rubbish-and-drains}
\hlchaptermodality{416}

\begin{chapteropening}
\textbf{By the end of this chapter:} \emph{I can complain about rubbish, a blocked drain and the smell.}

Complete the last lesson of chapter 416: I can complain about rubbish, a blocked drain and the smell.
\end{chapteropening}

\section[\texorpdfstring{\ur{کوڑا} (kūṛā)}{کوڑا (kūṛā)}]{\ur{کوڑا} (kūṛā) — rubbish}

\label{lesson:UR-C416-kura}



\begin{quote}
Before the new one: say the Urdu for a flavour, then the Urdu for to serve (food).
\end{quote}

\subsection*{You'll want to know: \ur{کوڑا}}
\textbf{\ur{کوڑا}} — \emph{kūṛā} — \textquotedblleft{}rubbish\textquotedblright{}.

\begin{quote}
\ur{گلی} \ur{میں} \ur{کوڑا} \ur{نہیں} \ur{پھینکنا} \ur{چاہیے۔}
\end{quote}

\begin{quote}
\emph{galī meṅ kūṛā nahīṅ pheṅknā chāhiye.}
\end{quote}

\begin{quote}
You shouldn't throw rubbish in the lane.
\end{quote}

Masculine.

\begin{cousinweb}
\textbf{\ur{کوڑا}} (\emph{kūṛā}) is inherited: CDIAL lists the Hindi cousin \emph{kūṛā}, \textquotedblleft{}a dust heap\textquotedblright{}, under Sanskrit \emph{kūṭa}, \textquotedblleft{}a heap\textquotedblright{} (CDIAL 3392).
\end{cousinweb}

\begin{grammarlens}[title={What you've built}]
One more word for getting around.
\end{grammarlens}

\subsection*{Your turn}
\begin{itemize}
\raggedright
  \item \emph{Say it:} \emph{kūṛā}
  \item \emph{Say it:} \emph{kūṛā}, once more
  \item \emph{Say it:} \emph{galī meṅ kūṛā nahīṅ pheṅknā chāhiye.} — You shouldn't throw rubbish in the lane.
  \item \emph{Recall:} say the Urdu for a flavour, then the Urdu for to serve (food), then say \emph{kūṛā} again
\end{itemize}

\begin{tcolorbox}[breakable,colback=teal!4,colframe=teal!35!black,title={Before you move on}]
What does \ur{کوڑا} mean? (\textquotedblleft{}Rubbish\textquotedblright{}.) Say it once more.
\end{tcolorbox}
\section[\texorpdfstring{\ur{نالی} (nālī)}{نالی (nālī)}]{\ur{نالی} (nālī) — a drain}

\label{lesson:UR-C416-nali}



\begin{quote}
Before the new one: say the Urdu for to serve (food), then the Urdu for rubbish.
\end{quote}

\subsection*{You'll want to know: \ur{نالی}}
\textbf{\ur{نالی}} — \emph{nālī} — \textquotedblleft{}a drain\textquotedblright{}.

\begin{quote}
\ur{بارش} \ur{کے} \ur{بعد} \ur{نالی} \ur{بھر} \ur{گئی۔}
\end{quote}

\begin{quote}
\emph{bārish ke bād nālī bhar gaī.}
\end{quote}

\begin{quote}
After the rain the drain filled up.
\end{quote}

Feminine.

\begin{cousinweb}
\textbf{\ur{نالی}} (\emph{nālī}) is built on \emph{nāl}, \textquotedblleft{}a hollow stalk, a tube\textquotedblright{}, which CDIAL lists under Sanskrit \emph{nāḍī}, \textquotedblleft{}a tube\textquotedblright{} (CDIAL 7047), with the ending \emph{-ī}.
\end{cousinweb}

\begin{grammarlens}[title={What you've built}]
One more word for getting around.
\end{grammarlens}

\subsection*{Your turn}
\begin{itemize}
\raggedright
  \item \emph{Say it:} \emph{nālī}
  \item \emph{Say it:} \emph{nālī}, once more
  \item \emph{Say it:} \emph{bārish ke bād nālī bhar gaī.} — After the rain the drain filled up.
  \item \emph{Recall:} say the Urdu for to serve (food), then the Urdu for rubbish, then say \emph{nālī} again
\end{itemize}

\begin{tcolorbox}[breakable,colback=teal!4,colframe=teal!35!black,title={Before you move on}]
What does \ur{نالی} mean? (\textquotedblleft{}A drain\textquotedblright{}.) Say it once more.
\end{tcolorbox}
\section[\texorpdfstring{\ur{گٹر} (gaṭar)}{گٹر (gaṭar)}]{\ur{گٹر} (gaṭar) — a sewer; a covered drain}

\label{lesson:UR-C416-gatar}



\begin{quote}
Before the new one: say the Urdu for rubbish, then the Urdu for a drain.
\end{quote}

\subsection*{You'll want to know: \ur{گٹر}}
\textbf{\ur{گٹر}} — \emph{gaṭar} — \textquotedblleft{}a sewer; a covered drain\textquotedblright{}.

\begin{quote}
\ur{سڑک} \ur{کے} \ur{نیچے} \ur{ایک} \ur{بڑا} \ur{گٹر} \ur{ہے۔}
\end{quote}

\begin{quote}
\emph{saṛak ke nīche ek baṛā gaṭar hai.}
\end{quote}

\begin{quote}
There's a big sewer under the road.
\end{quote}

Masculine.

\begin{cousinweb}
\textbf{\ur{گٹر}} (\emph{gaṭar}) is borrowed from English \emph{gutter}; in Urdu it names a sewer.
\end{cousinweb}

\begin{grammarlens}[title={What you've built}]
One more word for getting around.
\end{grammarlens}

\subsection*{Your turn}
\begin{itemize}
\raggedright
  \item \emph{Say it:} \emph{gaṭar}
  \item \emph{Say it:} \emph{gaṭar}, once more
  \item \emph{Say it:} \emph{saṛak ke nīche ek baṛā gaṭar hai.} — There's a big sewer under the road.
  \item \emph{Recall:} say the Urdu for rubbish, then the Urdu for a drain, then say \emph{gaṭar} again
\end{itemize}

\begin{tcolorbox}[breakable,colback=teal!4,colframe=teal!35!black,title={Before you move on}]
What does \ur{گٹر} mean? (\textquotedblleft{}A sewer; a covered drain\textquotedblright{}.) Say it once more.
\end{tcolorbox}
\section[\texorpdfstring{\ur{بدبو} (badbū)}{بدبو (badbū)}]{\ur{بدبو} (badbū) — a stink; a bad smell}

\label{lesson:UR-C416-badbu}



\begin{quote}
Before the new one: say the Urdu for a drain, then the Urdu for a sewer.
\end{quote}

\subsection*{You'll want to know: \ur{بدبو}}
\textbf{\ur{بدبو}} — \emph{badbū} — \textquotedblleft{}a stink; a bad smell\textquotedblright{}.

\begin{quote}
\ur{گٹر} \ur{سے} \ur{بدبو} \ur{آتی} \ur{ہے۔}
\end{quote}

\begin{quote}
\emph{gaṭar se badbū ātī hai.}
\end{quote}

\begin{quote}
The sewer stinks.
\end{quote}

Feminine.

\begin{cousinweb}
\textbf{\ur{بدبو}} (\emph{badbū}) is built from Persian \emph{bad}, \textquotedblleft{}bad\textquotedblright{}, and \emph{bū}, \textquotedblleft{}a smell\textquotedblright{}: it is usually traced to Persian \emph{badbū}.
\end{cousinweb}

\begin{grammarlens}[title={What you've built}]
One more word for getting around.
\end{grammarlens}

\subsection*{Your turn}
\begin{itemize}
\raggedright
  \item \emph{Say it:} \emph{badbū}
  \item \emph{Say it:} \emph{badbū}, once more
  \item \emph{Say it:} \emph{gaṭar se badbū ātī hai.} — The sewer stinks.
  \item \emph{Recall:} say the Urdu for a drain, then the Urdu for a sewer, then say \emph{badbū} again
\end{itemize}

\begin{tcolorbox}[breakable,colback=teal!4,colframe=teal!35!black,title={Before you move on}]
What does \ur{بدبو} mean? (\textquotedblleft{}A stink; a bad smell\textquotedblright{}.) Say it once more.
\end{tcolorbox}
\section[\texorpdfstring{\ur{پھیلانا} (phailānā)}{پھیلانا (phailānā)}]{\ur{پھیلانا} (phailānā) — to spread (something) around; to strew}

\label{lesson:UR-C416-phailana}



\begin{quote}
Before the new one: say the Urdu for a sewer, then the Urdu for a stink.
\end{quote}

\subsection*{You'll want to know: \ur{پھیلانا}}
\textbf{\ur{پھیلانا}} — \emph{phailānā} — \textquotedblleft{}to spread (something) around; to strew\textquotedblright{}.

\begin{quote}
\ur{کچھ} \ur{لوگ} \ur{گلی} \ur{میں} \ur{کوڑا} \ur{پھیلاتے} \ur{ہیں۔}
\end{quote}

\begin{quote}
\emph{kuchh log galī meṅ kūṛā phailāte haiṅ.}
\end{quote}

\begin{quote}
Some people strew rubbish about the lane.
\end{quote}

Past, with \ur{نے} (he or she …): \ur{اس} \ur{نے} \ur{میز} \ur{پر} \ur{نقشہ} \ur{پھیلایا}, \emph{us ne mez par naqshā phailāyā}.

\begin{cousinweb}
\textbf{\ur{پھیلانا}} (\emph{phailānā}) is built on \textbf{\ur{پھیلنا}}, \textquotedblleft{}to spread\textquotedblright{}: CDIAL lists both, as Hindi \emph{phailnā} and \emph{phailānā}, under Sanskrit \emph{prathita}, \textquotedblleft{}spread out\textquotedblright{} (CDIAL 8651). This is the form for spreading something else.
\end{cousinweb}

\begin{grammarlens}[title={What you've built}]
One more word for getting around.
\end{grammarlens}

\subsection*{Your turn}
\begin{itemize}
\raggedright
  \item \emph{Say it:} \emph{phailānā}
  \item \emph{Say it:} \emph{phailānā}, once more
  \item \emph{Say it:} \emph{kuchh log galī meṅ kūṛā phailāte haiṅ.} — Some people strew rubbish about the lane.
  \item \emph{Recall:} say the Urdu for a sewer, then the Urdu for a stink, then say \emph{phailānā} again
\end{itemize}

\begin{tcolorbox}[breakable,colback=teal!4,colframe=teal!35!black,title={Before you move on}]
What does \ur{پھیلانا} mean? (\textquotedblleft{}To spread (something) around; to strew\textquotedblright{}.) Say it once more.
\end{tcolorbox}
